%======================================================================
\section{Conclusions}
\label{sec:conclusions}

\subsection{What was known}
Pointwise, $\delta_k=\tau_k-\gamma_k$ is bounded, with $\limsup|\delta_k|\le0.6331$ from the explicit bound of Bellotti and Wong \cite{BellottiWong2025}, and $\limsup|\delta_k|\le0.6365$ from Trudgian's asymptotic bound \cite{Trudgian2010}. On average, $s_k$ is asymptotically Gaussian with variance $\frac1{2\pi^2}\log\log T$, by Selberg's formulae for the Gram displacement proved by Korolev \cite{Korolev2012,Korolev2011}, so $\langle\delta_k^2\rangle\sim2\log\log T/\log^2T$: the clock is close to the zeros in height but misplaces almost every zero by an unbounded number of positions. Selberg's approximation \cite{Selberg1946a,Fujii1974} replaces $S(t)$ by short prime sums in $t$-mean, and Fran\c{c}a--LeClair and LeClair \cite{FrancaLeClair2015,LeClair2021,FrancaLeClair2024} write the equation of an individual zero with an Euler phase and a growing number of primes, without an error estimate over the zeros.

\subsection{What this paper adds}
\begin{sloppypar}
\begin{enumerate}[leftmargin=1.6em]
\item \emph{An exact dictionary between clock-type models and the zeros} (Theorems~\CGref{thm:identity_general} and \CGref{thm:identity_g}). For any $C^1$ model $F$ and any continuous $\mathfrak g$, sums of $\mathfrak g$-increments of $k-\frac12-F(\gamma_k)$ over the ordinates are exactly $dF$-integrals of $\mathfrak g(N-F)$ up to two boundary terms; for powers this gives the even moments with the constant $\frac1{12}$ per zero, for exponentials it transfers the signed law of $S(t)$ to $s_k$ (Corollary~\CGref{cor:signed_law}).
\item \emph{The prime information needed by the clock} (Theorems~\CGref{thm:main} and \CGref{thm:variable_cutoff}). The corrected clock equation $\NO(t)+S_X(t)=k-\frac12$ holds at the true ordinates with mean-square ordinal defect $\frac1{2\pi^2}\log\frac{\log T}{\log X}+O(1+\sqrt{\log(\log T/\log X)})$, also with a cutoff varying at every index; the defect is bounded in mean square if and only if the primes reach a fixed power of the height.
\item \emph{An exact finite-height account} (Theorem~\CGref{thm:accounting}). The gain of the one-step clock is a diagonal prime response minus a cross term, and its residual differs from the defect by a transport term and an explicit nonnegative second-order term.
\item \emph{A computable clock that approaches the zeros} (Theorems~\CGref{thm:computable_clock}, \CGref{thm:second_step} and \CGref{thm:uniform_iteration}). With $X=T^{1/\sqrt{\log\log T}}$ the one-step clock has mean-square error $(\log\log\log T+O(\sqrt{\log\log\log T}))/\log^2T$ in height, against $(2+o(1))\log\log T/\log^2T$ for the bare clock; the $m$-th Newton step has error $\frac{H_m}{4\pi^2}+O(1+\sqrt{H_m})$ in units of zeros with absolute constants, and after $\log^*T+O(1)$ steps the error is $O(1)$ mean spacings, using no ordinate. The proof needs an exact moment recursion, a multiplicity bound for the corrected points of every level (Theorem~\CGref{thm:multiplicity_all}) and a cutoff law that absorbs the multiplicity.
\item \emph{The pointwise envelope} (Proposition~\CGref{prop:envelope_equivalence}, Theorem~\CGref{thm:envelope}). $\limsup|\tau_k-\gamma_k|=2\pi\limsup|S(t)|/\log t$ exactly, no slowly growing prime correction changes it, and Backlund's method with the profile generated from Bourgain's exponent pair gives $\limsup|\tau_k-\gamma_k|<0.60831$.
\end{enumerate}
\end{sloppypar}

\subsection{Open problems}
\begin{question}\label{q:constant}
For fixed $c\in(0,\frac12]$, do the bounded means of $u_k(T^c)^2$ converge, and to $(2\pi^2)^{-1}\log(1/c)$ plus which constant?
\end{question}

\begin{question}\label{q:fixed_c}
Does $\langle r_k(T^c)^2\rangle_{\IT}$ stay bounded for fixed $c$, that is, does the transport cancel against the defect? Is the error of the corrected clock after $\log^*T+O(1)$ steps bounded in units of zeros for \emph{almost every} index, and not only in mean square? Can Theorems~\CGref{thm:main} and \CGref{thm:computable_clock} be extended to $c>\frac12$?
\end{question}

\begin{question}\label{q:explicit}
Can the absolute constants $C_\star$, $H^*$ and $B_\star$ of Theorem~\CGref{thm:uniform_iteration}, and those of Fujii's and Montgomery--Vaughan's bounds on which they rest, be made explicit, so that the cutoff law \CGeqref{eq:cutoff_law} becomes a formula that can be evaluated at a given height?
\end{question}

Theorem~\CGref{thm:identity_general} applies to any counting staircase and any $C^1$ model, in particular to the zeros of a primitive Dirichlet $L$-function; an analogue of Theorem~\CGref{thm:main} would additionally need Selberg's approximation for $L(s,\chi)$ \cite{Selberg1946b,Fujii1974} and uniformity in the conductor.

\section*{Declaration of competing interest}
The author declares no competing interests.

\section*{Data availability}
No data were used: every statement is proved. The two programs of Appendix~\CGref{app:programs} evaluate the integral \CGeqref{eq:envelope_formula} for piecewise linear exponent profiles with rational vertices; they and their outputs are in the public GitHub repository \url{https://github.com/vcasana1/LawCasanya}, directory \texttt{work/paper1/}. They used Python 3 and mpmath 1.3.0.

\section*{Declaration of generative AI and AI-assisted technologies in the writing process}
During the preparation of this work the author used an AI assistant (Claude, Anthropic) to edit the text, check proofs and write programs. The author reviewed and edited the content and takes full responsibility for it.

%======================================================================
